% ── Abstract and keywords ────────────────────────────────────
\begin{abstract}
Running a quantum circuit fault-tolerantly on a surface-code processor reduces
to two coupled tasks: placing logical qubits on a two-dimensional lattice and
scheduling the lattice-surgery operations that execute its gates, each
non-Clifford gate consuming a magic state. Unlike the near-term (NISQ) setting,
the bottleneck is not $\SWAP{}$ insertion but the routing corridors---free lanes
of lattice that must stay open so interacting patches can be joined on demand. We
present a \emph{Gaussian potential-field} heuristic for the placement
sub-problem: each candidate tile is scored by superposing attractive Gaussian
fields---pulling $T$-heavy qubits toward magic states and tightly interacting
qubits toward their partners---with repulsive fields that preserve routing space,
and a \emph{safe-passage} test keeps every mapping routable. We implement the
method in an open-source C++20 OpenQASM compiler and evaluate it against a
randomized-restart baseline and the WISQ compiler. On \todo{$N$} circuits the
field placement reduces routing steps by \todo{$X$}\% on average over the
randomized baseline and is competitive with WISQ---completing several instances
on which WISQ times out---at a fraction of the compilation time.
\end{abstract}

\begin{IEEEkeywords}
qubit mapping, qubit routing, fault-tolerant compilation, surface code,
lattice surgery, magic states, quantum circuit compilation, heuristic placement
\end{IEEEkeywords}
